\section{Sign tests and the last moving site}
\label{sign:section}

We finish the exclusion of a finite critical wall.
\Cref{model:angle-budget-activated} supplies both the angle budget and
the polynomial bounds of \Cref{cor:critical-polynomial-bounds}.
In particular, with
\(d_k=e^{-k}\) and horizontal means normalized to have mass one,
\begin{equation}\label{sign:polynomial-input}
 \mu(k)\le C(1+k),\qquad
 d_k\sup_x\bigl(|f|+\sqrt{|q|}\bigr)
       (x+\ii(H-d_k))\le C(1+k)^2.
\end{equation}
Real symmetry gives the same estimates at the lower wall.  All constants
below may depend on the fixed billiard and on fixed tile enlargements.
They never depend on the large parameters \(K\) and \(L\).

Fix throughout
\begin{equation}\label{sign:alpha}
 \alpha=\frac1{100}.
\end{equation}
The purpose of the first part of the section is to find a deterministic
shape row on which almost every point has either very small or very large
\(|f|\).  Both alternatives suppress the motion of that site in the odd
Jacobi system, but by different identities.  The final argument applies
this suppression to the last site to cross the tested row.

\subsection{A discrete occupation estimate}

We first isolate the probabilistic estimate needed when the row moments
are small.  It does not require a bound on the starting value of the
submartingale.

\begin{lemma}[Discrete Tanaka estimate]\label{lem:discrete-tanaka}
Let \((X_m)_{0\le m\le M}\) be an integrable real submartingale, and suppose
\[
 V=\sum_{m=0}^{M-1}\E(X_{m+1}-X_m)^2<\infty.
\]
For \(0<h<R\), put
\[
 \mathcal C_h=
 \sum_{m=0}^{M-1}\E\left[
 \ind_{\{|X_m|\le h\}}
 \min\{(X_{m+1}-X_m)^2,h^2\}\right].
\]
There is an absolute constant \(C\) such that
\begin{equation}\label{sign:tanaka-bound}
 \mathcal C_h\le
 C h(h+R+\sqrt V)
 +C M h^2\min\left\{1,\frac{V}{(R-h)^2}\right\}.
\end{equation}
\end{lemma}

\begin{proof}
Use the increasing convex function
\[
 \Phi(x)=\sqrt{h^2+x^2}+x.
\]
Its derivative is between \(0\) and \(2\), it is nonnegative, and
\(\Phi(x)\le h+2x^+\).  Moreover,
\(\Phi''(x)\ge c/h\) on \([-2h,2h]\).  Integrating this second derivative
along the first \(\min\{|y-x|,h\}\) of the segment from \(x\) to \(y\)
proves
\begin{equation}\label{sign:ramp-remainder}
 \Phi(y)-\Phi(x)-\Phi'(x)(y-x)
 \ge \frac c h\ind_{\{|x|\le h\}}\min\{(y-x)^2,h^2\}.
\end{equation}
This includes a displacement that leaves \([-2h,2h]\), because the
remaining convex remainder is nonnegative.

Let \(T\) be the first index at which \(X_m>R\), capped at \(M\).
If \(X_0>R\), set \(T=0\).  Sum \eqref{sign:ramp-remainder} only on
\(m<T\).  Conditional expectation removes the linear term with the correct
sign, since \(\Phi'\ge0\) and the conditional drift is nonnegative.
The resulting sum is at most
\(C h\E[\Phi(X_T)-\Phi(X_0)]\).  Paths starting above \(R\) contribute
zero.  On the others the terminal positive part is at most
\[
 R+\max_{m<M}|X_{m+1}-X_m|,
\]
including the crossing increment.  Cauchy--Schwarz therefore bounds the
pre-\(T\) cost by \(C h(h+R+\sqrt V)\).

Write
\[
 b_m=\E[X_{m+1}-X_m\mid\mathcal F_m]\ge0,\qquad
 Z_n=\sum_{m<n}\ind_{\{T\le m\}}
       (X_{m+1}-X_m-b_m).
\]
This is a square-integrable martingale starting at zero, with terminal
second moment at most \(V\).  If a chain which starts above \(R\), or
crosses \(R\), subsequently returns to \([-h,h]\), then
\(\sup_n|Z_n|\ge R-h\); accumulated drift cannot assist this downward
move.  The martingale maximal inequality bounds the probability of any
such return by \(C\min\{1,V/(R-h)^2\}\).  After \(T\), the entire cost is
at most \(Mh^2\) on that event.  Adding the two estimates proves
\eqref{sign:tanaka-bound}.  In particular, no conditional estimate given
the crossing event has been used.
\end{proof}

\subsection{The large-moment alternative}

\begin{lemma}\label{sign:large-moment}
Suppose that, along a sequence \(K\to\infty\), some
\(j\in[0.9K,2.1K]\) satisfies \(\mu(j)\ge K^\alpha\).  There is a
deterministic \(k\in[0.8K,2.2K]\) for which
\begin{equation}\label{sign:row-small-exception}
 \frac1{2\pi}\left|\left\{x\in[0,2\pi]:
 \left|\log|f(x+\ii(H-e^{-k}))|\right|
           \le K^{\alpha/2}\right\}\right|\longrightarrow0.
\end{equation}
\end{lemma}

\begin{proof}
Choose a small fixed \(c>0\), with \(c<1/4\), and maximize
\(\mu(l)e^{-c|l-j|}\) over \(l\ge k_0\).  The maximum exists by
\eqref{sign:polynomial-input}.  A maximizer \(k\) satisfies
\[
 |k-j|=O(\log K),\qquad B:=\mu(k)\ge K^\alpha,
 \qquad \mu(l)\le B e^{c|l-k|}.
\]
For a uniform random center \(x\), use the normalized profiles
\begin{equation}\label{sign:profile-normalization}
 F_{K,x}(Z)=\frac{d_k}{B}f(x+\ii H+d_k Z),\qquad
 Q_{K,x}(Z)=\frac{d_k^2}{B^2}q(x+\ii H+d_k Z).
\end{equation}
Their domains exhaust the lower half-plane. At a depth in the fixed
upper collar, \(0<D\le e^{k-k_0}\), the moment envelope gives
\begin{equation}\label{sign:profile-moments}
 \E\bigl(|F_{K,x}(X-\ii D)|^2+|Q_{K,x}(X-\ii D)|\bigr)
 \le C D^{-2}e^{2c|\log D|}.
\end{equation}
The bound is uniform in the horizontal coordinate \(X\).  Submean
estimates on fixed disks give tightness in the topology of locally
uniform convergence.  More explicitly, take a countable compact
exhaustion, view the restrictions in the product of the corresponding
continuous-function spaces, and put a zero placeholder on compact sets
not yet contained in a domain.  Tightness and diagonal extraction give
consistent holomorphic limits.  Uniformity of the original center makes
each limiting law stationary under all real translations.

The scaled equation is
\[
 F_{K,x}''=B^2Q_{K,x}F_{K,x}.
\]
Consequently each limit \((F,Q)\) is pure: \(FQ=0\), so at least one of
the two holomorphic functions vanishes identically.  The limit law has no
atom at \((0,0)\).  To check precisely the extra hypothesis needed for
\Cref{thm:no-trivial-profile}, choose
\(1/2<\eta<1-c\). In a fixed enlarged cone, subdivide the coarser
depths \(1\le D\le\delta_b/d_k\) into dyadic bands. The choice of
\(\delta_b\) in \Cref{sec:critical} keeps even the fixed tile
enlargements within \(d_kD\le e^{-k_0}\), the range of
\eqref{sign:profile-moments}. Each band contains only a bounded
number of tiles of size comparable to its depth.  By
\eqref{sign:profile-moments} and submean, the second moments of
\[
 D^\eta\bigl(|F_{K,x}|+\sqrt{|Q_{K,x}|}\bigr)
\]
on these tiles have a summable bound
\(C D^{-2(1-c-\eta)}\).  Thus the required cone bound has a tight random
constant.  Positive mass at the trivial pair would allow selection of
centers converging to that pair while this constant stayed bounded:
first discard a set of probability smaller than half the alleged atom,
then use successively smaller compact neighborhoods of zero.  The
moment envelope has coarser exponent \(c<1/2\) and finer exponent
\(c>0\), so this selection contradicts
\Cref{thm:no-trivial-profile}.

There is one more degeneracy to exclude.  With probability one, a
nonzero limiting \(Q\) is not real on the entire line \(\Im Z=-1\).
Otherwise condition on the event that it is real there and nonzero;
this event is translation invariant.  Schwarz reflection across that
line extends \(Q\) to an entire function. For each fixed \(D\ge1\),
\eqref{sign:profile-moments} holds for all sufficiently large \(k\),
so Fatou's lemma gives
\[
 \E|Q(X-\ii D)|\le C D^{-2(1-c)}.
\]
Conditioning increases the constant only by the reciprocal probability
of the event.  Reflection supplies the same bound in the reflected
half-plane.  Every point is represented either at depth at least one or
by its reflection at such a depth; hence the conditional mean modulus
is bounded on the entire plane.  Cauchy's derivative estimate on circles
of radius \(R\), followed by expectation, gives
\(\E|Q'(z)|\le C/R\).  Letting \(R\to\infty\) makes \(Q\) constant
almost surely; the decay at large depth makes that constant zero.  This
contradicts the event on which we conditioned.

We now examine any deterministic convergent sequence of profiles with
one of the allowed limits.  If \(F\not\equiv0\), then, away from its
discrete zeros on a fixed row segment,
\[
 \log|f(x+\ii H+d_kZ)|=k+\log B+\log|F_{K,x}(Z)|.
\]
The right side exceeds \(K^{\alpha/2}\) uniformly on compact subsets
away from those zeros.  Arbitrarily small neighborhoods of the zeros
therefore contain all the exceptional points in the limit.

Suppose instead that \(Q\not\equiv0\).  Since its imaginary part on
\(\Im Z=-1\) is a nonzero real-analytic function, its zeros on compact
row intervals are discrete.  Apart from arbitrarily small
neighborhoods of these points, cover the row by finitely many short
intervals with complex neighborhoods on which consistent branches
\(p_K=\sqrt{Q_{K,x}}\) satisfy
\[
 \Re p_K\ge c_1>0
\]
and have uniformly bounded derivatives.  In the scaled variable the
original function \(g(Z)=f(x+\ii H+d_kZ)\) satisfies
\(g''=B^2Q_{K,x}g\).  A basis on each such interval has the form
\begin{equation}\label{sign:wkb-basis}
 g_\pm(u)=Q_{K,x}(u-\ii)^{-1/4}
  \exp\left(\pm B\int_{u_0}^u p_K(v-\ii)\,\dd v\right)
  (1+O(B^{-1})).
\end{equation}
Here and below the error is uniform on the chosen closed interval.
For completeness, the expressions without the final factor solve the
same equation with an added uniformly bounded potential; their
Wronskian has modulus comparable to \(B\).  Variation of constants,
from the left for the growing expression and from the right for the
decaying expression, gives Volterra equations for the relative
factors.  After division by the leading expressions the nonconstant
kernel contains a factor
\[
 \exp\left(-2B\int_v^u p_K(r-\ii)\,\dd r\right),\qquad v\le u,
\]
whose modulus is at most \(e^{-2Bc_1(u-v)}\); the remaining kernel
has size \(O(B^{-1})\).  Successive approximation, or the integral
Gronwall inequality, therefore gives relative error \(O(B^{-1})\).
The two solutions are independent, since their modulus ratio changes
by an exponential factor along the interval.  This proves
\eqref{sign:wkb-basis}.

Write \(g=A_+g_++A_-g_-\), with coefficients allowed to depend on
\(K\).  If both coefficients are nonzero, the difference of the two
leading logarithmic magnitudes has derivative at least \(2c_1B-O(1)\).
The set on which that difference has modulus at most \(\sqrt B\)
has length \(O(B^{-1/2})\).  Outside it one term dominates the other by
at least \(e^{\sqrt B+O(1)}\).  Each individual leading logarithmic
magnitude is strictly monotone with slope of absolute value at least
\(c_1B-O(1)\), so it spends at most
\(O((K^{\alpha/2}+1)/B)=o(1)\) length within
\(K^{\alpha/2}+1\) of zero.  The same conclusion holds when one
coefficient vanishes.  Thus the fraction violating
\eqref{sign:row-small-exception} on every fixed scaled interval tends
to zero.

The countable product of continuous-function spaces used above is
Polish. By the Skorokhod representation theorem
\cite[Chapter~10, Theorem~4, p.~239]{Pollard2002}, any weakly
convergent subsequence of profile laws admits an almost-sure coupling.
The deterministic criterion therefore applies almost surely.
The local bad fraction is bounded by one, so its expectation tends to
zero.  Translation averaging identifies this expectation with the
full-period bad fraction on the selected row.  Every subsequence has
the same conclusion, proving the lemma.
\end{proof}

\subsection{Amplitude control when the row moments are small}

Assume for the remainder of this subsection that
\begin{equation}\label{sign:low-moments}
 \mu(v)\le K^\alpha\qquad (0.9K\le v\le2.1K).
\end{equation}
For a shape point \(z\), write
\[
 \delta(z)=H-|\Im z|,\qquad
 a(z)=\log|f(z)|,\qquad
 \mathcal A(z)=\log\bigl(|f(z)|+\delta(z)|f'(z)|\bigr).
\]
The function \(\mathcal A\) is finite and continuous: a common zero of
\(f\) and \(f'\) would imply \(f\equiv0\) by the equation
\(f''=qf\).  In contrast, \(a\) has logarithmic singularities at the
isolated zeros of \(f\).  The next estimate keeps these two issues
separate.

\begin{lemma}[Local amplitude and logarithmic deficit]
\label{sign:amplitude}
Let \(z_0=x+\ii(H-d)\), \(d<H/4\), and put
\[
 H_b=1+\sup_{B(z_0,d/2)}\delta\sqrt{|q|}.
\]
On \(B(z_0,d/8)\), the oscillation of \(\mathcal A\) is at most
\(CH_b\).  Moreover, the normalized area \(L^2\) norm of
\(\mathcal A-a\) on that disk, and its normalized length \(L^2\) norm
on the horizontal diameter, are at most \(CH_b\).  The same estimates
hold for fixed smaller disks and for horizontal segments of length
comparable to \(d\) with fixed interior margins.
\end{lemma}

\begin{proof}
Scale by \(d/8\), writing
\[
 F(\zeta)=f(z_0+(d/8)\zeta),\qquad
 Q(\zeta)=(d/8)^2q(z_0+(d/8)\zeta).
\]
These functions are holomorphic on \(|\zeta|<4\), and
\(|Q|^{1/2}\le H_b/4\).  The nonvanishing vector
\(Y=(F,F'/H_b)\) satisfies
\[
 Y'=\begin{pmatrix}0&H_b\\ Q/H_b&0\end{pmatrix}Y.
\]
The matrix norm is at most \(CH_b\).  Gronwall along a straight
segment, used in both directions, bounds the change of
\(\log\|Y\|\) by \(CH_b\) on fixed interior disks.  Since physical
depth there is comparable to \(d\), replacing \(\log\|Y\|\) by
\(\mathcal A\) changes the oscillation by at most
\(C\log(2+H_b)\).  In particular,
\begin{equation}\label{sign:amplitude-oscillation}
 |\mathcal A(z)-\mathcal A(z_0)|\le CH_b
       \qquad (|z-z_0|\le3d/8).
\end{equation}

Set \(A_0=\mathcal A(z_0)\) and \(G=e^{-A_0}F\).
Cauchy's estimate on a sufficiently small fixed circle shows that
there is a point \(|\zeta_*|\le1/16\) with
\(\log|G(\zeta_*)|\ge-C\): indeed a fixed multiple of
\(\sup_{|\zeta|\le1/16}|F(\zeta)|\) bounds
\(|F(0)|+8|F'(0)|=e^{A_0}\).  On \(|\zeta|\le3\),
\eqref{sign:amplitude-oscillation} gives \(\log|G|\le CH_b\).
Choose \(M= C_0H_b\) large enough and set
\(U=M-\log|G|\ge0\).  Then \(U(\zeta_*)\le CH_b\).

Jensen's formula on a circle centered at \(\zeta_*\), of radius in
\([2,9/4]\) avoiding boundary zeros, bounds by \(CH_b\) the number of
zeros in \(B(\zeta_*,7/4)\), with multiplicities.  Choose a further
radius in \([3/2,7/4]\) avoiding boundary zeros.  The boundary mean
of \(U\) is at most \(U(\zeta_*)\), by submean for \(\log|G|\).
Poisson--Jensen expresses \(U\) inside this latter circle as the
harmonic extension of that nonnegative boundary data plus the positive
disk Green functions of the zeros.  On \(|\zeta|\le1\), the harmonic
term is bounded by \(CH_b\).  Each Green function is bounded there by
\[
 C+\log^-|\zeta-\zeta_j|,
\]
which has uniformly bounded \(L^2\) norm on the unit disk and on its
horizontal diameter.  For the segment assertion use
\(\log^-|t+\ii b-\zeta_j|\le\log^-|t-\Re\zeta_j|\) and the
integrability of \(|\log|t||^2\).  There are at most \(CH_b\) zeros,
so the triangle inequality bounds the required \(L^2\) norms of
\(U\) by \(CH_b\).  Finally,
\[
 \mathcal A-a=(\mathcal A-A_0-M)+U,
\]
and \eqref{sign:amplitude-oscillation} bounds the other summand.
Fixed smaller disks and comparable segments follow by the same fixed
ratios, or by a finite covering.
\end{proof}

The tile quantities in \Cref{sign:amplitude} obey
\begin{equation}\label{sign:tile-L2}
 \E_x H_b(x,d_v)^2\le CK^{2\alpha},\qquad
 \E_x(\mathcal A-a)^2(x+\ii(H-d_v))\le CK^{2\alpha},
\end{equation}
whenever the fixed enlargement stays within the band
\eqref{sign:low-moments}.  To see the first assertion, submean for
\(|q|\), followed by averaging in the center, gives
\[
 \E_x H_b^2\le C\left(1+
 d_v\int_{d_v/4}^{7d_v/4}
 \E_x|q(x+\ii(H-t))|\,\dd t\right)\le CK^{2\alpha}.
\]
The second follows by integrating the segment estimate of
\Cref{sign:amplitude} in its center.  Larger fixed tiles or apertures
cost only a fixed covering factor.  The lower row follows by symmetry.

Run ordinary Brownian motion in the shape strip from a uniform point
of its real period.  At the first hit of either row of depth \(e^{-v}\),
write \(a(v)\) for the value of \(a\).  These variables, sampled on
any increasing finite grid, form an integrable submartingale.  Indeed
apply subharmonic comparison in the next substrip to
\(\max\{\log|f|,-T\}\) and then let \(T\to\infty\); row logarithms
are integrable, including at isolated zeros.  Translation invariance
makes the horizontal coordinate uniform at each sample, and symmetry
makes the two sides equiprobable.

\begin{lemma}[Logarithmic increments]\label{sign:increments}
Under \eqref{sign:low-moments}, uniformly for
\(1.1K\le v\le1.9K\) and \(2\le s\le4\),
\begin{equation}\label{sign:increment-L2}
 \E|a(v+s)-a(v)|^2\le CK^{2\alpha}.
\end{equation}
\end{lemma}

\begin{proof}
Start on the upper row of depth \(d=e^{-v}\) and stop on first
hitting either row of depth \(d'=de^{-s}\).  Let \(U\) be the
unwrapped horizontal displacement.  The distance to the upper target
line is \(d-d'\asymp d\).  Ignoring the lower target line gives
half-plane exit with Cauchy density
\[
 \frac{d-d'}{\pi((d-d')^2+u^2)}.
\]
Consequently
\begin{equation}\label{sign:jump-tails}
 \Prob(\text{same side},\ |U|>Rd)\le C/R\quad(R\ge1),
 \qquad \Prob(\text{opposite side})\le Cd.
\end{equation}
The second estimate is also the elementary linear hitting probability
between two horizontal lines.  Conditional on displacement and on
both sides, the starting horizontal point is still uniform modulo
\(2\pi\); the terminal point is its translate, hence uniform as well.
This follows directly by translating the Brownian path and its start.

For a same-side displacement \(u\) with \(R=|u|/d\le e^{K/20}\),
join the endpoints by a prescribed path going first to depth
\(D=d+|u|\), then horizontally by \(u\), and then to depth \(d'\).
Subdivide the vertical pieces in fixed small relative steps.  The
horizontal piece requires only a bounded number of depth-sized tiles.
The whole chain has at most \(C(1+\log(1+R))\) tiles.  All their fixed
enlargements lie in the low-moment band for large \(K\).  At fixed
\(u\), their centers are fixed translates of the uniform starting
point.  By \Cref{sign:amplitude}, \eqref{sign:tile-L2}, and Minkowski,
\[
 \|\mathcal A(\text{end})-\mathcal A(\text{start})\|_{L^2(dx/(2\pi))}
 \le CK^\alpha(1+\log(1+R)).
\]
The square of the last factor has bounded expectation by
\eqref{sign:jump-tails}.

For all remaining transitions, compare amplitudes by a chain from
each endpoint to a fixed real collar, connecting there modulo period.
The global bound \eqref{sign:polynomial-input} makes the sum of the
tile costs at most \(C(1+K)^3\): there are \(O(K)\) levels and the
cost at level \(j\) is \(O((1+j)^2)\).  Thus the amplitude difference
is pointwise polynomially bounded even when an endpoint is near a
zero of \(f\).  The remaining transitions have probability at most
\(C(e^{-K/20}+e^{-v})\), by \eqref{sign:jump-tails}.  They contribute
\(o(1)\) to its second moment.  We have proved
\[
 \E|\mathcal A(\text{end})-\mathcal A(\text{start})|^2
       \le CK^{2\alpha}.
\]
Finally subtract \(\mathcal A-a\) at both endpoints.  Their uniform
horizontal marginal laws and \eqref{sign:tile-L2} bound each squared
error by \(CK^{2\alpha}\).  This proves
\eqref{sign:increment-L2}.  Only the amplitude difference was bounded
pointwise on rare transitions; no pointwise bound for the logarithmic
singularities is needed.
\end{proof}
\subsection{From a curvature witness to an occupation cost}

Set
\begin{equation}\label{sign:h0}
 h_0=K^{6\alpha}.
\end{equation}
Apply \Cref{lem:discrete-tanaka} to step-\(s\) grids in
\([1.1K,1.9K]\), including the last increment whose initial point is
in the interval.  The number of increments is \(O(K)\), and
\Cref{sign:increments} gives \(V\le CK^{1+2\alpha}\).  Choose
\(R=K^{1/2+10\alpha}\).  Then
\[
 h_0(h_0+R+\sqrt V)=O(K^{1/2+16\alpha}),\qquad
 K h_0^2\frac{V}{(R-h_0)^2}=O(K^{1-6\alpha}).
\]
Both quantities are \(o(K/\log K)\) for \eqref{sign:alpha}.
Integrate the grid estimate over its initial offset in an interval of
length \(s\), and then over \(2\le s\le4\).  Every starting level is
counted once.  We obtain
\begin{equation}\label{sign:occupation-cost}
 \int_{1.1K}^{1.9K}\int_2^4
 \E\left[\ind_{\{|a(v)|\le h_0\}}
 \min\{|a(v+s)-a(v)|^2,h_0^2\}\right]\,\dd s\,\dd v
       =o(K/\log K).
\end{equation}
The estimates at the bounded number of terminal increments are covered
by the margins in \eqref{sign:low-moments}.

\begin{lemma}[Witnesses force logarithmic oscillation]
\label{sign:witness-cost}
Under \eqref{sign:low-moments}, the fraction of pairs
\((x,k)\in[0,2\pi]\times[1.2K,1.8K]\) for which
\[
 |a(x+\ii(H-e^{-k}))|\le h_0/4
\]
tends to zero.
\end{lemma}

\begin{proof}
Choose the aperture and all fixed enlargements large enough for
\Cref{cone:no-quiet-witness}, with the polynomial exponent in
\eqref{sign:polynomial-input}.  That proposition supplies fixed
constants \(L_0,C_1,C_2,c_*>0\).  At every target \((x,k)\), there
is a witness \(z\) with
\begin{equation}\label{sign:quiet-witness-window}
 |\Re z-x|\le L_0\delta(z),\qquad
 k-C_1\log k\le-\log\delta(z)\le k+C_2,
 \qquad \delta(z)(|f(z)|+\sqrt{|q(z)|})\ge c_*.
\end{equation}
Increasing the fixed margins allows all disks and paths below to stay
inside this enlarged cone window.

For a uniformly chosen target \((x,k)\) in the stated central band,
\eqref{sign:tile-L2} and Markov's inequality show, with probability
tending to one, that
\begin{equation}\label{sign:typical-target}
 (\mathcal A-a)(x+\ii(H-e^{-k}))\le K^{2\alpha},
 \qquad \delta\sqrt{|q|}\le K^{2\alpha}
\end{equation}
throughout every required enlarged conic tile.  For the second
assertion, cover the window by a bounded number of tiles in each unit
log-depth band and use the union bound over \(O(\log K)\) bands.
The failure probability is \(O(K^{-2\alpha}\log K)=o(1)\).

Call such a target typical.  If it is also moderate, meaning
\(|a|\le h_0/4\), then \Cref{sign:amplitude} along tile chains of
length \(O(1+\log K)\) gives
\begin{equation}\label{sign:cone-amplitude}
 |\mathcal A|\le h_0/3
\end{equation}
on the needed cone cells, for all large \(K\).  Indeed the initial
error and total variation are
\(O(K^{2\alpha}(1+\log K))=o(h_0)\).  On this entire window,
\(\delta|f|\le\exp(-1.1K+h_0/3)\) for large \(K\), so the
first field cannot provide \eqref{sign:quiet-witness-window}.
At a witness point \(z_1\) we therefore have
\begin{equation}\label{sign:q-witness}
 \delta_1\sqrt{|q(z_1)|}\ge c_*/2,
       \qquad \delta_1=\delta(z_1).
\end{equation}

Let \(\mathcal D_1\) be a disk centered at \(z_1\) of a sufficiently
small fixed radius times \(\delta_1\).  There is a fixed \(c_2>0\)
such that, for every real number \(z\),
\begin{equation}\label{sign:log-oscillation}
 \frac1{|\mathcal D_1|}\int_{\mathcal D_1}
       |a-z|\,\dd\mathrm{area}\ge c_2.
\end{equation}
Here \(|\mathcal D_1|\) denotes its area.  To prove uniformity,
suppose the assertion failed along a
sequence.  Scale the disks to a fixed disk and multiply \(f\) by
\(e^{-z}\).  The resulting holomorphic functions \(g_n\) satisfy
\(\log|g_n|\to0\) in \(L^1\).  Submean bounds \(|g_n|\) uniformly
on compact subdisks, so a subsequence converges holomorphically.
The limit is not zero, since uniform convergence to zero on a
subdisk would contradict the \(L^1\) convergence of the logarithms.
Off its zeros its modulus is one, again by that convergence.  Thus
it is a nonzero constant.  At the center,
\(g_n''/g_n\to0\), contradicting \eqref{sign:q-witness} in scaled
coordinates.  This proves \eqref{sign:log-oscillation} without an
upper bound for the witness curvature independent of \(K\).

For \(|z|\le h_0/2\), the lower bound persists after restricting to
\(|a|\le h_0\) and truncating \(|a-z|\) at \(h_0\).  Indeed,
\eqref{sign:cone-amplitude} and \Cref{sign:amplitude} give
\[
 \frac1{|\mathcal D_1|}\int_{\mathcal D_1}(\mathcal A-a)^2
                 \,\dd\mathrm{area}\le CK^{4\alpha}.
\]
On \(|a|>h_0\), the upper bound \(a\le\mathcal A\le h_0/3\)
forces \(a<-h_0\), so \(\mathcal A-a\ge2h_0/3\).
The discarded \(L^1\) mass of \(|a-z|\) is at most
\(CK^{4\alpha}/h_0=o(1)\), by Cauchy--Schwarz or the elementary
tail estimate \(\E[e\ind_{e>t}]\le\E e^2/t\).
On the retained set, \(|a-z|\le3h_0/2\), and truncation at \(h_0\)
retains at least two thirds of it.  Cauchy--Schwarz now yields
\begin{equation}\label{sign:truncated-witness}
 \frac1{|\mathcal D_1|}\int_{\mathcal D_1}\ind_{\{|a|\le h_0\}}
          \min\{|a-z|^2,h_0^2\}\,\dd\mathrm{area}
       \ge c_3>0\qquad (|z|\le h_0/2).
\end{equation}

Take a second disk \(\mathcal D_2\), with the same horizontal
center, centered at depth \(e^{-3}\delta_1\), and with radius a
small fixed multiple of that depth.  It lies in the same enlarged
cone window.  By \eqref{sign:cone-amplitude} and the disk deficit
estimate, all but \(o(1)\) of its normalized area has
\(|a|\le h_0/2\).  For every starting point of \(\mathcal D_1\),
the transitions to the upper row at an additional log depth
\(s\in[2,4]\), integrated in \(s\), dominate a fixed positive
multiple of normalized area on \(\mathcal D_2\).  Here the radii
were chosen small enough that all height increments lie in \([2,4]\)
with room.  The upper-exit strip Poisson kernel has density at least
\(c/\delta_1\) for these displacements, since both horizontal
displacement and vertical distance to the target row are
\(O(\delta_1)\), and the latter is comparable to \(\delta_1\).
Also normalized area at depth comparable to \(\delta_1\) is
comparable to \(\dd x\,\dd v/\delta_1\), with only fixed factors
for the second disk.

For a starting point \(x'+\ii(H-e^{-v})\), let
\[
 C(x',v)=\int_2^4\E_{x',v}\left[
 \ind_{\{|a(x',v)|\le h_0\}}
 \min\{|a(v+s)-a(x',v)|^2,h_0^2\}\right]\,\dd s,
\]
where the conditional expectation denotes ordinary shape Brownian
exit to the two next rows, and \(a(x',v)\) abbreviates the starting
value.  Integrate the kernel lower bound over \(\mathcal D_1\),
then over the terminal points of \(\mathcal D_2\) with
\(|a|\le h_0/2\).  Formula \eqref{sign:truncated-witness}, applied
with each such terminal value as \(z\), shows that every typical
moderate target forces
\begin{equation}\label{sign:target-cost}
 1\le C\int_{\mathcal W(x,k)}C(x',v)
                     \frac{\dd x'\,\dd v}{e^{-v}},
\end{equation}
where \(\mathcal W(x,k)\) is the full enlarged cone window.
Replacing the witness disk by this window is legitimate because the
integrand is nonnegative; it also avoids any measurable-choice issue
for the witness.

Integrate \eqref{sign:target-cost} over all typical moderate targets.
For a fixed source \((x',v)\), the horizontal measure of target
centers whose window contains it is \(O(e^{-v})\); this cancels the
factor \(e^v\).  Their allowed \(k\)-range has length
\(O(\log K)\).  Fubini therefore bounds their unnormalized measure
by
\[
 C\log K\int_{1.1K}^{1.9K}\int_0^{2\pi}C(x',v)\,\dd x'\,\dd v.
\]
Uniform horizontal law and the upper-side probability \(1/2\)
bound the last integral by a fixed multiple of
\eqref{sign:occupation-cost}.  The result is \(o(K)\).  Atypical
targets already have fraction \(o(1)\), proving the lemma.
\end{proof}

\begin{theorem}[Deterministic row sign test]\label{thm:row-sign}
For every sufficiently large \(K\), there is a deterministic
\(k=k(K)\in[0.8K,2.2K]\) such that, on each shape row of depth
\(e^{-k}\),
\begin{equation}\label{sign:row-test}
 |\log|f||>K^{\alpha/2}
\end{equation}
outside a set of relative horizontal measure tending to zero as
\(K\to\infty\).
\end{theorem}

\begin{proof}
If \(\mu(j)\ge K^\alpha\) for some
\(j\in[0.9K,2.1K]\), use \Cref{sign:large-moment}.  Otherwise
\eqref{sign:low-moments} holds, and
\Cref{sign:witness-cost} provides a deterministic row in
\([1.2K,1.8K]\) with exceptional fraction tending to zero for the
stronger threshold \(h_0/4\).  Since
\(K^{\alpha/2}=o(h_0)\), it also gives \eqref{sign:row-test}.
Real symmetry gives the lower-row conclusion from the upper one.
The sequential proofs imply the stated conclusion as \(K\to\infty\),
without restricting to either alternative along a fixed subsequence.
\end{proof}

\subsection{Robustness of the sign test}

\begin{proposition}[A robust sign disk]\label{prop:robust-sign}
For the deterministic level in \Cref{thm:row-sign}, put
\begin{equation}\label{sign:h1-radius}
 h_1=\tfrac12K^{\alpha/2},\qquad
 r_{\mathrm{rob}}=e^{-k-h_1/4}.
\end{equation}
Outside a set of relative horizontal measure tending to zero on
either tested row, each point \(z\) has
\[
 |\log|f(\zeta)||\ge h_1
       \qquad (|\zeta-z|\le r_{\mathrm{rob}}),
\]
with constant sign on that disk.  In the odd-slice Brownian
experiment, for each fixed site the probability that its first hit
of the tested shape depth occurs at a point without this property
tends to zero.
\end{proposition}

\begin{proof}
Write \(d=e^{-k}\).  By \eqref{sign:polynomial-input},
\(|q|\le C d^{-2}K^4\) throughout a fixed depth-sized neighborhood
of either tested row.  Zeros of \(f\) in this neighborhood are
separated by at least \(c dK^{-2}\).  Indeed, scale around a zero
by \(r=c dK^{-2}\), choosing \(c\) small.  The scaled potential
has arbitrarily small supremum on a fixed disk, and the solution
normalized to have derivative one and value zero at the center
satisfies
\[
 F(\zeta)=\zeta+
       \int_0^\zeta(\zeta-u)Q(u)F(u)\,\dd u.
\]
The integral equation gives \(|F(\zeta)-\zeta|\le C\|Q\|_\infty
|\zeta|^3\) on a smaller fixed disk.  Thus the zero is the only
one there.  This also reproves simplicity in this setting.

Packing disks of radius comparable to \(dK^{-2}\) in a periodic
band of width \(O(d)\) bounds its number of zeros by
\(C d^{-1}K^4\).  Delete from the tested row the shadows of disks
of radius
\[
 r_{\mathrm{exc}}=d e^{-h_1/10}
\]
about all these zeros.  Their total relative length is at most
\(CK^4e^{-h_1/10}=o(1)\).  Zeros outside this band are too far away
to affect the following smaller disks.

At a remaining point, and throughout its disk of radius
\(r_{\mathrm{rob}}\),
\begin{equation}\label{sign:log-derivative}
 \left|\frac{f'}f\right|\le C/r_{\mathrm{exc}}.
\end{equation}
Here is the local zero test proving this bound.  If
\(\lambda=f'(z)/f(z)\) has modulus much larger than
\(r_{\mathrm{exc}}^{-1}\), rescale by \(1/\lambda\) and divide by
\(f(z)\).  On \(|u|\le2\) the resulting equation has arbitrarily
small potential, because
\(\sup|q|/|\lambda|^2\le CK^4e^{-h_1/5}\), and its initial data
are \(F(0)=F'(0)=1\).  The integral equation makes \(F\) uniformly
close to \(1+u\).  Rouche's theorem gives a zero near \(u=-1\),
hence within \(2/|\lambda|\) of \(z\).  That contradicts the
distance to the deleted zero disks.  Since
\(r_{\mathrm{rob}}/r_{\mathrm{exc}}=e^{-3h_1/20}\), the same
argument applies throughout the robustness disk.

The disk is zero-free, so integrating \eqref{sign:log-derivative}
along segments bounds the change of \(\log|f|\) by
\(Ce^{-3h_1/20}=o(1)\).  The original test has threshold \(2h_1\),
which proves the assertion after discarding its exceptional set.
For the Brownian statement, \Cref{odd:time-change} dominates the
unconditional first-hit law of each site by a fixed multiple of
ordinary shape Brownian motion from uniform horizontal start.
The latter has uniform horizontal distribution on each tested row.
The bad first-hit probability is therefore \(o(1)\), including
paths on which the site is stopped before reaching that row.
\end{proof}
\subsection{The last moving site}

Return to the odd periodic phase slice of
\Cref{lem:odd-matrices,thm:stopped-geometry}.  The odd integer \(N\)
and the real parameter \(t_0=2\pi/N\) are fixed.  Brownian time in
the phase strip is denoted by \(u\), so derivatives \(s_j'\) always
mean derivatives with respect to the complex phase coordinate \(y\),
not derivatives in time.  Retain the notation
\[
 a_j=\log|f(s_j)|,\qquad
 k_j=-\log(H-|\Im s_j|),\qquad
 K_j(u)=\sup_{v\le u}k_j(v),\qquad
 w_{\max}=\max_j|w_j|.
\]
The stopping time \(\tau_L\) is the first time one of the histories
\(K_j\) reaches \(L\).  By \Cref{thm:boundary-contact}, it occurs
almost surely before phase exit for all large \(L\).

\begin{lemma}[A deterministic history gap]\label{sign:history-gap}
There is a fixed finite menu of exponents \(0<\gamma_m<1\) such
that for each large \(L\), one deterministic choice
\(K=L^{\gamma_m}\) has the following event of probability greater
than \(0.9\): every terminal history \(K_j(\tau_L)\) is either
at most \(K^{\alpha/8}\) or at least \(3K\).  At least one history
is in the second class.
\end{lemma}

\begin{proof}
Choose a fixed integer \(M>10N\).  Choose positive exponents
\(\gamma_1<\cdots<\gamma_M<1\), small enough and separated so that
\(\gamma_m<\gamma_{m+1}\alpha/8\).  Then the intervals
\[
 (L^{\gamma_m\alpha/8},3L^{\gamma_m}),\qquad 1\le m\le M,
\]
are pairwise disjoint and lie below \(L\), for all large \(L\).
Each path has at most \(N\) failed menu choices: a single terminal
history belongs to at most one of the intervals.  Averaging failure
probabilities over the menu gives a deterministic choice with
failure probability at most \(N/M<0.1\).  One history equals \(L\)
at \(\tau_L\), so the large class is nonempty.  Because the menu
is fixed and finite, all subsequent \(o(1)\) estimates can be made
uniform over its choices.
\end{proof}

Fix this deterministic \(K\), and set
\begin{equation}\label{sign:h2}
 u_0=K^{\alpha/8},\qquad h_2=K^{\alpha/4}=u_0^2.
\end{equation}
On the event of \Cref{sign:history-gap}, denote the shallow indices
by \(I\) and their nonempty complement by \(J\).  These sets depend
on the terminal path.  In the estimates below, however, every
candidate set is fixed before a stopping argument is applied.

\begin{lemma}[Two derivative bounds]\label{sign:derivative-bounds}
With exceptions of probability tending to zero, the following hold
through \(\tau_L\).  For every site,
\begin{equation}\label{sign:site-suppression}
 |s_i'|\le e^{o(h_2)-|a_i|}w_{\max}.
\end{equation}
For each fixed nonempty candidate \(J\), put \(I=J^c\) and stop
additionally when any history in \(I\) reaches \(2u_0\).  Up to
that stop,
\begin{equation}\label{sign:deep-controls-all}
 w_{\max}\le e^{Ch_2}\max_{j\in J}|s_j'|.
\end{equation}
The exceptional probability tends to zero uniformly over the finite
family of candidates.  In particular, both bounds hold throughout
\([0,\tau_L]\) on the final classification event with its actual
sets \(I,J\), apart from an event of probability tending to zero.
\end{lemma}

\begin{proof}
Choose the arbitrarily small fixed exponent in
\Cref{thm:stopped-geometry} using the smallest exponent in the menu.
Its geometry cost is then \(o(u_0)\), uniformly through \(\tau_L\).
This controls logarithms of the sines and lengths, and hence all
entries of \(M_0\), at cost \(o(h_2)\).
The identities \eqref{odd:w-definition} and
\eqref{odd:jacobi-angle} read
\begin{equation}\label{sign:two-identities}
 w_i=c\sin d_i^0\,f_i^{-2}s_i',\qquad
       2f_i s_i'=(M_0w)_i.
\end{equation}
If \(a_i<0\), the first gives
\(|s_i'|\le e^{o(h_2)+2a_i}w_{\max}\), which implies
\eqref{sign:site-suppression}.  If \(a_i\ge0\), the second gives
that inequality directly.

We spell out the uniform inverse estimate needed for the second
bound.  Fix \(I\) and a subset \(F\subset J\).  Retain the sites
\(I\) with their actual curvatures, retain \(F\) with curvature
zero, and delete all other sites.  This is one of the generically
invertible modified matrices of \Cref{lem:odd-matrices}; it has at
least one deletion or flattening, since \(J\ne\varnothing\).
Stop at the first shallow history \(2u_0\), and at the geometry
cuts just selected.  The upper bound for \(a_i\) on truthful sites
is \(O(u_0)\) by \eqref{sign:polynomial-input}.  Every entry and
the determinant of this fixed matrix thus have upper logarithmic
bound \(Cu_0\).  At a real starting point its determinant is
nonzero, with a uniform lower bound on the compact real period.
Apply the stopped logarithmic-deficit estimate from
\Cref{odd:log-deficit}, with lower threshold \(u_0^{3/2}\).
Except with probability \(O(u_0^{-1/2})\), the inverse determinant
costs at most \(e^{u_0^{3/2}}\).  The adjugate costs \(e^{Cu_0}\),
so
\begin{equation}\label{sign:modified-inverse}
 \|B_{I,F}^{-1}\|
       \le \exp(O(u_0^{3/2}))=e^{o(h_2)}.
\end{equation}
Empty determinants are one.  This is also the specialization of
\Cref{odd:modified-inverses} with upper threshold \(O(u_0)\) and
lower threshold \(u_0^{3/2}\).  Take the finite union over all
fixed pairs \((I,F)\) before using any time-dependent set.

At a given time put
\[
 F=\{j\in J:a_j<-h_2\},\qquad D=J\setminus F.
\]
In \(\mathcal H w=0\), eliminate the block \(I\cup F\) by
\eqref{sign:modified-inverse}.  The curvatures removed on \(F\)
satisfy
\[
 |\mathcal V_j|=
 \left|\frac{2f_j^3}{c\sin d_j^0}\right|
       \le e^{-3h_2+o(h_2)}.
\]
Their diagonal perturbation is absorbed by the inverse bound.
The coupling entries to \(D\) cost only \(e^{o(h_2)}\).
Consequently \(w_{\max}\le e^{o(h_2)}\max_{j\in D}|w_j|\).
If \(D\) is empty the same absorption gives \(w=0\), so the
desired inequality still holds.  Otherwise \(a_j\ge-h_2\) on
\(D\), and the first identity in \eqref{sign:two-identities}
gives
\[
 |w_j|\le e^{2h_2+o(h_2)}|s_j'|.
\]
This proves \eqref{sign:deep-controls-all}.  All possible flat sets
were already included in the finite union; no assertion about their
switching times is needed.  On the final shallow event the additional
history cut is after \(\tau_L\), which proves the last assertion.
\end{proof}

\begin{theorem}[The last-site contradiction]\label{thm:last-site}
The nonisotropic finite critical-wall case is impossible.
\end{theorem}

\begin{proof}
Use \Cref{thm:row-sign} for the deterministic \(K\) just chosen,
and let \(k\in[0.8K,2.2K]\) be its deterministic tested level.
Write \(d=e^{-k}\), and retain
\(h_1=K^{\alpha/2}/2\) and
\(r_{\mathrm{rob}}=d e^{-h_1/4}\) from
\Cref{prop:robust-sign}.  Notice the strict separation
\begin{equation}\label{sign:scale-separation}
 1\ll u_0\ll u_0^{3/2}\ll h_2\ll h_1\ll K.
\end{equation}

For each site \(j\), let \(\sigma_j\) be its first hit of shape
depth \(d\), and put \(\sigma_j=\infty\) if it never hits that
depth before phase exit.  Define the remaining intrinsic clock
\[
 Q_j=\int_{\sigma_j\wedge\tau_L}^{\tau_L}|s_j'(y_u)|^2\,\dd u.
\]
The shape Brownian time-change extension in
\Cref{odd:time-change} and the strong Markov property give
\begin{equation}\label{sign:remaining-clock}
 \Prob\{Q_j>d^2e^{\sqrt{h_2}}\}
       \le C e^{-\sqrt{h_2}/2}.
\end{equation}
Indeed, after first hit the point is distance \(d\) from its
nearest wall.  Exit from the shape strip occurs no later than first
hit of that wall in the Brownian extension.  For a one-dimensional
Brownian motion started at distance \(d\) from a line, the reflection
principle gives
\(\Prob(T_{\mathrm{line}}>t)\le C d/\sqrt t\).
Substitute \(t=d^2e^{\sqrt{h_2}}\).  If the tested row has not been
reached at \(\tau_L\), then \(Q_j=0\).  Thus
\eqref{sign:remaining-clock} is unconditional and does not depend
on later membership of \(j\) in a classification set.

For every fixed nonempty subset \(J\), define
\begin{equation}\label{sign:last-hit-stop}
 T_J=\max_{j\in J}(\sigma_j\wedge\tau_L).
\end{equation}
This is a stopping time, as a finite maximum of stopping times, and
\(T_J\le\tau_L\).  On the event that \(J\) is the actual deep
class, every one of its sites reaches history \(3K\).  Since
\(k<3K\), every \(\sigma_j\) is then strictly before
\(\tau_L\), and \(T_J\) is an actual last first hit.
On the good geometry, inverse, and clock events, the bound
\eqref{sign:deep-controls-all} yields, for the actual classification,
\begin{equation}\label{sign:w-clock}
 \int_{T_J}^{\tau_L}w_{\max}^2\,\dd u
 \le e^{Ch_2}\sum_{j\in J}
       \int_{\sigma_j}^{\tau_L}|s_j'|^2\,\dd u
 \le d^2e^{C'h_2}.
\end{equation}
The constants absorb the fixed number of sites and the squaring of
the exponential factor in \eqref{sign:deep-controls-all}.

We next give the displacement estimate with fixed candidates; this
is necessary because the terminal classification and last-hit order
are not known at time zero.  Fix \(J\ne\varnothing\) and \(i\in J\).
At time \(T_J\), center a disk of radius \(r_{\mathrm{rob}}\) at
the \(\mathcal F_{T_J}\)-measurable point \(s_i(y_{T_J})\).
Let \(\eta_{J,i}\) be its first exit time after \(T_J\), capped at
\(\tau_L\).  Consider the event that \(J\) is the actual deep
class and \(i\) is an actual last site, so \(\sigma_i=T_J\).
Exclude the first-hit failures from \Cref{prop:robust-sign}, whose
unconditional probability is \(o(1)\) after the finite union over
sites.  On the event under consideration, \(|a_i|\ge h_1\) up to
\(\eta_{J,i}\).  By \eqref{sign:site-suppression} and
\eqref{sign:w-clock},
\begin{equation}\label{sign:suppressed-clock}
 \int_{T_J}^{\eta_{J,i}}|s_i'|^2\,\dd u
       \le d^2\exp(C''h_2-2h_1)=:q_K.
\end{equation}
Choose the fixed constant \(C''\) with enough slack to absorb all
the deterministic \(o(h_2)\) costs.

Apply the maximal inequality to the actual complex displacement
martingale
\(s_i(y_u)-s_i(y_{T_J})\), stopped at \(\eta_{J,i}\) and also
when its scalar quadratic time after \(T_J\) reaches \(2q_K\).
Its expected squared terminal displacement is at most \(Cq_K\).
One may first localize inside compact substrips and then pass to
the limit.  On the good event in \eqref{sign:suppressed-clock},
the additional clock cut does not precede \(\eta_{J,i}\).
Therefore the probability of disk exit on that good event is at
most
\begin{equation}\label{sign:escape-probability}
 C\frac{q_K}{r_{\mathrm{rob}}^2}
 \le C\exp(C''h_2-\tfrac32h_1)=o(1).
\end{equation}
This is an unconditional martingale estimate for each fixed
\((J,i)\).  The classification event, last-hit identity, geometry
bounds, and robustness are intersected with it only afterward.
Ties among last sites cause no problem: choose any tied index and
include it in the same finite union.

\begin{figure}[t]
\centering
\begin{tikzpicture}[>=Stealth,x=1.65cm,y=1.15cm]
\draw[thick] (-2.7,2.4)--(2.7,2.4) node[right,font=\small] {shape wall};
\draw[dotted] (-2.7,2.08)--(2.7,2.08);
\node[font=\small,anchor=west] at (2.8,2.08) {depth $e^{-3K}$};
\draw[dotted] (-2.7,.60)--(2.7,.60);
\node[font=\small,anchor=west] at (2.8,.60) {depth $e^{-k}$};
\fill[deepblue!8] (0,.60) circle (.33);
\draw[deepblue] (0,.60) circle (.33);
\fill[deepblue] (0,.60) circle (1.6pt);
\node[font=\small,anchor=east] at (-.46,.57) {$s_i(y_{T_J})$};
\draw[->,warmred,thick] (-.7,-.3) to[out=50,in=-110] (0,.60);
\draw[->,warmred,dashed,thick] (0,.60) to[out=65,in=-95] (.5,2.08);
\node[font=\small,align=left,anchor=west] at (.58,1.25)
 {Reaching the row nearer the wall\\requires leaving the disk.};
\draw[->,deepblue] (0,.60)--(.29,.75);
\node[font=\small,align=center] at (0,-.79)
 {Robust radius $e^{-k-h_1/4}$\qquad quadratic-time cost
 $\le e^{-2k+Ch_2-2h_1}$};
\end{tikzpicture}
\caption{The last-site contradiction, drawn without metric scale.
For a fixed candidate set, \(T_J\) is the last of its capped first hits
of the tested row. The last site's robust sign suppresses its subsequent
quadratic time so strongly that disk exit has vanishing probability.
Its prescribed final history would nevertheless require such an exit.}
\label{fig:last-site}
\end{figure}


Disk exit is nevertheless compulsory on the classification event.
An actual last site has depth \(d\) at \(T_J=\sigma_i\), and
subsequently attains depth at most \(e^{-3K}\) by \(\tau_L\).
It cannot have attained that depth earlier, since that would already
have crossed the tested row.  The depth function is
\(1\)-Lipschitz, so its displacement before \(\tau_L\) is at least
\[
 d-e^{-3K}\ge d/2>r_{\mathrm{rob}}
\]
for large \(K\).  Thus it must leave the disk shown in
\Cref{fig:last-site}.  The finite union of
\eqref{sign:escape-probability} over all candidates \(J,i\) makes
this event have probability \(o(1)\), whereas
\Cref{sign:history-gap} and all the excluded failure estimates give
probability at least \(0.9-o(1)\).  This contradiction proves the
theorem.
\end{proof}

\subsection{Completion of the rigidity theorem}

\begin{proof}[Proof of \Cref{thm:main}]
The physical caustic-collar hypothesis gives an analytic boundary by
\Cref{thm:analytic}, and then the even joint parametrization and
nonzero area coefficient of \Cref{thm:joint-parametrization}.
The continuation and no-patch theorems produce the maximal shape
strip \(D_H\), with the natural-boundary conclusion of
\Cref{cor:shape-natural-boundary} when \(H<\infty\).

If the meromorphic normal \(a=e^{\ii\phi}\) has a zero or pole
inside this strip, \Cref{thm:isotropic-classification} already
identifies the boundary as an ellipse.  In the remaining case,
\Cref{thm:shape-lift} and \Cref{lem:f-zeros} give the holomorphic
shape fields and the intrinsic impact lift used above.  Constant
\(f\) gives a circle.  For nonconstant \(f\),
\Cref{thm:infinite-growth-exclusion} excludes both \(H=\infty\)
and infinite upper growth order at a finite wall.  At a finite wall,
\Cref{prop:subcritical-exclusion} excludes upper growth order below
one, and \Cref{thm:finite-superlinear-exclusion} excludes every
finite upper growth order above one.  The only remaining possibility
is a finite critical wall.  The critical bootstrap,
\Cref{thm:sequence-exclusion}, and the odd-slice construction place
it under the hypotheses of \Cref{thm:last-site}, which excludes it
as well.

Only an elliptic boundary, including a circle, remains.  Convexity
then identifies the domain with the interior of that ellipse.
After translation, it is therefore described by a positive definite
quadratic form, as asserted.
\end{proof}
